% Evans 12.5, printed pp.642--647; English PDF pp.659--664.
% Chinese start: printed p.456 / scan p.474.
\section{Nonexistence of Solutions}
\label{sec:12-5}
\subsection{Negative energy and the concavity method}
Consider a classical solution of \eqref{eq:12-semi} with compactly
supported smooth data and sufficient regularity for the following
identities. Finite propagation bounds its support on every finite
strip. This section allows a negative potential.
\begin{theorem}[Source \S12.5.1, Theorem 1]
Suppose $f(0)=0$, $F'=f$, $F(0)=0$, and for some $\lambda>2$,
\[
 zf(z)\le\lambda F(z)\quad(z\in\R).
\]
If the conserved energy $E(0)<0$, no such solution can persist for
all $t\ge0$. In particular, the concavity argument forces breakdown
in a finite time.
\end{theorem}
\begin{expandedproof}[obtaining a decreasing concave quantity]
Put $I(t)=\frac12\int u^2$. Differentiating twice and integrating
by parts gives
\[
 I'=\int uu_t,\qquad I''=\int(u_t^2-|Du|^2-uf(u)).
\]
Set $\alpha=(\lambda-2)/4>0$ and $\beta=-\lambda E(0)>0$.
Energy conservation and the structural assumption yield
\[
 I''\ge(2+2\alpha)\int u_t^2+2\alpha\int|Du|^2+\beta,
 \qquad (I')^2\le2I\int u_t^2.
\]
Therefore $II''-(1+\alpha)(I')^2\ge\beta I$.
The function $I$ never vanishes: if $u=0$ on a time slice, its
energy there would be $\frac12\int u_t^2\ge0$.
For $J=I^{-\alpha}>0$,
\begin{equation}\label{eq:12-concavity}
 J''=-\alpha I^{-\alpha-2}
       \{II''-(1+\alpha)(I')^2\}
       \le-\alpha\beta I^{-\alpha-1}<0.
\end{equation}
If a global solution existed, $I''\ge\beta$ would make $I'$ positive
at some finite $t_0$. Then $J'(t_0)<0$, and concavity implies
$J(t)\le J(t_0)+(t-t_0)J'(t_0)$ for $t\ge t_0$.
The right side becomes zero in finite time, contradicting $J>0$.
No positivity assumption on the initial correlation $\int gh$ is
needed; $I''\ge\beta$ creates the required sign later.
\end{expandedproof}
\PDEAnnotation{负能量不需要初速度向外}{聚焦幂
$f(z)=-|z|^{p-1}z$ 满足 $zf=(p+1)F$。
取 $h=0,g=A\psi$，$A$ 足够大就使
$E=A^2\|D\psi\|_2^2/2-A^{p+1}\|\psi\|_{p+1}^{p+1}/(p+1)<0$。
证明不依赖 $I'(0)>0$。}

\subsection{Small-data blow-up below the three-dimensional threshold}
Now consider the different focusing equation
\begin{equation}\label{eq:12-source-power}
 \PDEWave u=|u|^p\quad\hbox{in }\R^3,
 \qquad \int g>0,\quad\int h>0,
 \quad\operatorname{spt}(g,h)\subset B(0,R_0).
\end{equation}
The data themselves need not be nonnegative. For a noninteger $p$
the statement concerns a classical solution with the regularity
needed for the representation and integrations, not an automatic
$C^\infty$ solution of a nonsmooth nonlinearity.
\begin{theorem}[Source \S12.5.2, Theorem 2]
If $1<p<1+\sqrt2$, the solution in \eqref{eq:12-source-power}
cannot exist classically for all positive time. The assumptions allow
arbitrarily small positive multiples of fixed admissible data.
The endpoint $p=1+\sqrt2$ is not established by this proof.
\end{theorem}
\begin{expandedproof}[the shell gain and a divergent time integral]
Assume a global solution and put $I(t)=\int u(x,t)dx$.
Finite propagation and integration of the PDE give
$I''=\int|u|^p$, hence $I,I'>0$. Holder on $B(0,R_0+t)$ gives
\begin{equation}\label{eq:12-I-diff}
 I''\ge c(1+t)^{-3(p-1)}I^p.
\end{equation}
Let $v$ be the free wave with the same data. Positivity of the
three-dimensional retarded kernel gives $u\ge v$, despite possible
sign changes of $v$. Huygens' principle puts $v$ in the shell
$A_t=\{\max(t-R_0,0)\le|x|\le t+R_0\}$, whose volume is at most
$C(1+t)^2$. Also $\int v=\int g+t\int h$.
Consequently
\[
 c(1+t)\le\int_{A_t}u
 \le\left(\int|u|^p\right)^{1/p}|A_t|^{1-1/p},
 \qquad I''\ge c(1+t)^{2-p}.
\]
Since $p<3$, two integrations and the positive initial values yield
$I(t)\ge c(1+t)^{4-p}$ (after adjusting the constant for small $t$).
This stronger lower bound comes from shell volume, not ball volume.

Choose $0<\varepsilon<p-1$ so that
$\mu=(p-1)^2+\varepsilon(4-p)<2$. Combining the lower bound with
\eqref{eq:12-I-diff} gives
$I''\ge c I^{1+\varepsilon}(1+t)^{-\mu}$.
Multiply by $I'>0$ and integrate. The product rule
\[
 (I^{2+\varepsilon}(1+t)^{-\mu})'
 =(2+\varepsilon)I^{1+\varepsilon}I'(1+t)^{-\mu}
       -\mu I^{2+\varepsilon}(1+t)^{-\mu-1}
\]
shows
$(I')^2\ge cI^{2+\varepsilon}(1+t)^{-\mu}-C$.
The positive term tends to infinity: its power lower bound has
exponent $2(4-p)-(p-1)^2>0$ on the stated range. Thus for large $t$,
\[
 I'\ge cI^{1+\varepsilon/2}(1+t)^{-\mu/2},\qquad
 (I^{-\varepsilon/2})'\le-c(1+t)^{-\mu/2}.
\]
Because $\mu/2<1$, the time integral on the right diverges.
It makes the positive left quantity negative at a finite time,
a contradiction. Scaling the initial data down changes constants
and the possible lifetime, but does not remove this contradiction.
\end{expandedproof}
\PDEAnnotation{两个临界指数不可混用}{这里 $1+\sqrt2$ 来自
$(p-1)^2<2$ 和球壳的 $t^2$ 体积增长；
三维退聚焦的能量临界 $p=5$ 来自缩放。
本节既没有证明 Strauss 临界端点，也没有讨论阈值以上的小数据全局理论。}
\PDEBodyDiagram{12.4}{Two different obstructions to a global solution}
{On the left, a positive concave quantity with negative derivative lies
below its tangent and must reach zero. The displayed curve
$J=1-t-t^2/4$ is only an exact concavity model, not a PDE solution.
On the right, $(p-1)^2<2$ leaves room for the positive $\varepsilon$
in the shell-based proof; the endpoint is excluded. Neither diagram
asserts a quantitative lifespan formula for the PDE.}{ch12-blowup}
